\documentclass[12pt,a4paper]{article}

% Encoding
\usepackage[utf8]{inputenc}
\usepackage[T1]{fontenc}

% Math
\usepackage{amsmath, amssymb}
% Images
\usepackage{graphicx}
% Links
\usepackage[colorlinks=true, urlcolor=blue, linkcolor=blue]{hyperref}

% Layout
\usepackage[a4paper, margin=1in]{geometry}
\setlength{\parindent}{0pt}

\title{\bfseries État de l'art \\ Explicabilité des CNN}
\author{SCHMITT Hugo}

\begin{document}

\maketitle

\section*{Contexte général : l'explicabilité ? ...}
Les CNNs, et NNs en général sont des boîtes noires...\\
Problèmes : confiance, robustesse, ... $\to$ champ de recherche : XAI


\section*{Taxonomie des méthodes d'explicabilité}

\subsection*{Méthodes par perturbation (indépendante du modèle)}
Comme exemple : l'Occlusion.\\
Principe : masquer une partie de l'entrée et observer la variation de sortie.\\
Intuition : importance d'une zone (selon taille du patch) = impact sur la prédiction.\\
Avantages : simple (compréhensible), interprétable directement, ..\\
Limites : coûteux (beaucoup de fp), dépend de la taille du patch.

\subsection*{Méthodes basées sur les gradients}
Comme exemple : Integrated Gradients.\\
Principe : comme un gradient de base, importance pour une entrée selon sa dérivée partielle, mais amélioré en intégrant les gradients entre une baseline (valeur de réference) et l'entrée...\\
Corrige donc les gradients bruités (reste bruyant), et la saturation des gradients.\\
Avantages : propriétés théoriques (axiomes).\\
Limites : dépend du choix de la baseline, parfois difficile à interpréter..

\subsection*{Méthodes basées sur les activations (CAM-like)}
Comme exemple : Grad-CAM.\\
Principe : utilise les gradients et les cartes d'activations convolutionnelles, fait ressortir les régions les plus importantes de l'entrée.\\
Avantages : très adapté aux CNNs, et donc interprétation visuelle intuitive.\\
Limites : résolution grossière?, dépend de la couche choisie.


\section*{Critères d'évaluation des méthodes}
Une explication n'a pas de vérité absolue nette. Pas de métrique absolue, seulement des approches cherchant à régler ce problème.\\
\textit{L'évaluation de l'explicabilité elle-même est un problème ouvert...}
\subsection*{Critères empirique}
Ces critères proviennent de l'expérimentation (d'où empirique), ne sont pas fondés.\\
Fidélité(faithfulness) : est-ce que l'explication reflète vraiment le modèle ?\\
Stabilité(robustness) : petites perturbations = mêmes explications ?\\
Localisation : précision spatiale.\\
Complexité computationnelle.

\subsection*{Axiomes}
Ils garantissent des propriétés mathématiques fondées (rigueur théorique), mais pas nécessairement une interprétation humaine pertinente.\\
Sensitivity...\\
Implementation Invariance...\\
Completeness...


\section*{Limites générales du domaine}

\subsection*{Instabilité des méthodes}
Deux méthodes =? deux explications différentes...

\subsection*{Sensibilité aux hyperparamètres}
baseline(IG),taille du patch(Occlusion),couche(Grad-CAM)

\subsection*{Problème fondamental}
Une explicabilité $\neq$ vérité du modèle.\\
Risque de fausse interprétation.


\section*{Tendances récentes}

\subsection*{Combinaison de méthodes (multi-perspective)}

\subsection*{Améliorations des méthodes existantes}

\subsection*{Méthodes globales / conceptuelles}

\end{document}